\chapter{What a period is not}

\textbf{Purpose:} Record a confident reading that is wrong, say exactly what it is reading instead,
and give the test that catches it. \textbf{Scope:}
\texttt{examples/\allowbreak{}game\_\allowbreak{}theory/\allowbreak{}6\_\allowbreak{}oracle/\allowbreak{}sturmian\_\allowbreak{}convergents.\allowbreak{}py},
\texttt{src/\allowbreak{}engine/\allowbreak{}python/\allowbreak{}measure/\allowbreak{}periodicity.\allowbreak{}py}

\section{The reading}

Take Wythoff's losing positions in order and write down the gaps between successive first
coordinates. Two values appear, 1 and 2, in ratio 1.62069 against \(\varphi = 1.61803\), and the word
they make is Sturmian. It has no period. That is not a measurement and no longer sequence would
change it: it follows from \(\varphi\) being irrational.

Hand that word to \texttt{sequence\_period} and it returns a number.

\begin{center}
\begin{tabular}{@{}lr@{}}
\toprule
window & reported \\
\midrule
16 & 8 \\
24 & 8 \\
32 & 13 \\
48 & 21 \\
64 & 21 \\
96 & 34 \\
128 & 55 \\
\bottomrule
\end{tabular}
\end{center}

Those are Fibonacci numbers. Every one of the seven rows clears its own shuffle floor, and over the
full table of nine slopes the margin runs 3.4 to 29 times that floor. No margin test refuses any of
them.

\section{What it is actually reading}

The mechanism is continued fractions and it is not mysterious.

The word is the gap sequence of a rotation by an irrational \(\alpha\). Such a word agrees with
itself shifted by \(q\) whenever \(q\alpha\) is close to an integer, and the \(q\) for which
\(q\alpha\) is closest to an integer are exactly the denominators of the convergents of \(\alpha\).
The Fibonacci numbers are the convergent denominators of \(\varphi\). So the detector is finding
real self-similarity, at the lags where real self-similarity is, and reporting it correctly. The
only thing wrong is the name on the output.

The prediction that follows is falsifiable and was tested. If the mechanism is continued fractions,
the reported number should be a convergent denominator of the slope for any irrational slope and not
only for \(\varphi\). Nine slopes, seven windows each:

\begin{center}
\begin{tabular}{@{}ll@{}}
\toprule
slope & reported, by window \\
\midrule
\(\varphi\) & 8, 8, 13, 21, 21, 34, 55 \\
\(\sqrt 2\) & 5, 12, 12, 17, 29, 29, 29 \\
\(e\) & 7, 7, 7, 7, 32, 39, 39 \\
\(\pi\) & 7, 7, 7, 7, 7, 7, 53 \\
\(\sqrt 3\) & 4, 11, 15, 15, 15, 41, 56 \\
\(\sqrt 5\) & 4, 4, 13, 17, 17, 17, 55 \\
\(\sqrt 7\) & 3, 11, 14, 17, 31, 48, 48 \\
\(\ln 2 + 1\) & 8, 10, 13, 13, 13, 13, 62 \\
\(\sqrt[3]{2} + 1\) & 4, 4, 15, 23, 27, 27, 50 \\
\bottomrule
\end{tabular}
\end{center}

Sixty one of the sixty three reported numbers are a convergent or intermediate denominator of their
own slope. The intermediate fractions are counted because they are also one-sided best
approximations and a detector reading agreement has no reason to prefer the two-sided ones.

The two exceptions are the two weakest rows in the table. Sorted by margin, they sit at rank 0 and
rank 8 of 63, and the lowest of all of them is one of the two. Where the mechanism does not
account for the answer, the detector was also at its least sure.

\(\pi\) is the row that tests the mechanism, since a fit alone would not predict it. Its continued
fraction has a 292 in it early. Its convergent denominators therefore jump from 7 to 106, and the
detector holds 7 across six consecutive windows before moving at all. A slope whose approximations improve slowly,
like \(\varphi\), moves at every step. The shape of the movement is the shape of the slope's
continued fraction.

\section{The reading is a capability read backwards}

Nothing here says the detector is broken. Handed a word built from a rotation, its answers across
widening windows enumerate the convergent denominators of the rotation's slope, which is to say it
recovers the continued fraction of a number it was never given. That belongs in the same list as an
image returning its own width and a protein backbone returning its bond lengths.

What it cannot be allowed to do is return that number under the word \emph{period} to a caller who
will read it as one. This work has one instrument reporting periods, the crystallography chapter
uses it on cell edges, and a reader who believes a high margin means a period believes something the
table above refutes.

\section{The test, and it costs nothing}

A true period is a property of the sequence. An approximation is a property of the sequence and the
window together. So widen the window and see whether the answer moves.

\begin{center}
\begin{tabular}{@{}lr@{}}
\toprule
& one answer at every window \\
\midrule
subtraction games with a real period, at 5 windows & 115 of 115 \\
aperiodic words, at 7 windows & 0 of 9 \\
\bottomrule
\end{tabular}
\end{center}

The separation is total on the rows tested, it needs no oracle, no null and no second instrument,
and it costs one extra call per window. It is in the instrument as
\texttt{measure.\allowbreak{}periodicity.\allowbreak{}stable\_\allowbreak{}period}, which returns
a period only where widening the window leaves the answer alone and returns nothing where it does
not. \texttt{sequence\_period} is unchanged, because this is a second question and not a correction
to that one. Any reading of a period anywhere in this work that was taken
at a single window is, on this evidence, unverified. It is not shown to be wrong. The crystallography
rows are safe because a lattice has a period for the reported number to be.

\section{What this does not reach}

The window sweep separates a period from a rotation. It has not been shown to separate a period from
anything else, and the aperiodic words tested are all Sturmian, which is one family. A sequence that
is aperiodic in some other way might well hold its answer across windows and defeat the test. Nine
slopes is nine, the same way ten domains is ten, and this chapter is subject to the workbook's own
rule that an eleventh instance moves the quantifier no distance.

The claim that holds is narrow and it is enough: a confident margin is not evidence of a period, one
worked mechanism explains what the confident margin was reading instead, and a test that separates
the two on every row tried is free.
